\subsection{流产后的心理调适与伴侣支持}

上一节讲了人工流产的术式与医学禁忌（禁欲 1 个月、高效避孕、月经恢复时间）。这一节专门讲常被略过的另一半——\textbf{心理与关系的恢复}。无论是有计划的终止妊娠，还是意外失去胎儿，流产后的情绪波动都是普遍且正常的。

\vspace{0.5em}
\noindent\textbf{常见的情绪反应}：

\begin{itemize}
	\item \textbf{情绪起伏是常态而非异常}：失落、内疚、愤怒、空虚、反复回想决策过程……这些反应可持续数周甚至数月；\textbf{激素水平的骤然回落}（妊娠相关激素在流产后数天至两周内大幅下降）会放大情绪波动——"身体也在经历一场分娩式的退潮"
	\item \textbf{两种常见的自我攻击}：一是"是不是我做了什么"（加班、吃了药、同房了）——绝大多数流产是\textbf{胚胎染色体异常}（早期流产的 50\%--60\%）等随机因素所致，与孕妇的日常行为几乎无关；二是"我不配难过，因为是我自己决定不要的"——\textbf{决定艰难与悲伤真实，两者并不矛盾}
	\item \textbf{需要专业帮助的信号}：情绪低落、失眠、自责持续超过 4--6 周且无缓解；出现自伤念头；或发展为抑郁发作——此时应就诊心理科/精神科，流产后抑郁是真实存在的临床问题，不是"想开点"能解决的
\end{itemize}

\vspace{0.5em}
\noindent\textbf{伴侣的角色}：

\begin{itemize}
	\item \textbf{双方哀悼节奏可能不同}：女性承受身体创伤与激素波动，情绪反应常更直接；男性一方则可能压抑情绪、假装"没事"来"扛住"——各哀各的、互不理解，是流产后伴侣疏远的常见剧本
	\item \textbf{有效的支持姿态}：陪伴而不催促（"你想说的时候我都在"）、不做评判（不说"反正还没准备好"）、\textbf{主动分担具体的恢复事务}（买卫生用品、复诊陪同、家务）
	\item \textbf{别急着"再来一次"}：双方都需要时间处理丧失，切莫把"尽快再怀一个"当作修复手段——那会把新的怀孕变成止痛药
\end{itemize}

\vspace{0.5em}
\noindent\textbf{身体恢复与再次备孕的节奏}：

\begin{itemize}
	\item \textbf{禁欲与避孕}：术后 1 个月内禁欲（详见上节）；之后若无生育计划，\textbf{立即恢复高效避孕}——流产后 2 周左右就可能恢复排卵，"月子期不会怀孕"是危险的误会
	\item \textbf{月经恢复}：多在 4--6 周；超过 6 周未至、或恢复后经量极少伴周期性腹痛，需复诊排除宫腔粘连
	\item \textbf{再次备孕时机}：一般建议\textbf{恢复 1--3 次正常月经后}再试孕，让内膜与心理都有缓冲；早期自然流产者若身心状态良好、月经已恢复，可与医生讨论提前；反复流产（$\geq$2 次）者应先做系统评估（染色体、免疫、解剖因素）再备孕
\end{itemize}

\begin{tcolorbox}[colback=pink!5!white,colframe=red!50!black,title=贴心小叮咛]
流产是一段\textbf{真实的丧失}，哪怕它发生在很早期。给身体一个月，给心情一段时间，给关系一点耐心——三者都恢复好了，再往前走不迟。如果一个月后你发现自己仍被情绪困住，请把这当作需要就医的事，而不是需要忍耐的事。
\end{tcolorbox}

